% ============================================================
%  CHAPTER 5 — EVALUATION FRAMEWORK
%  Fully rewritten 2026-08-28 -- see VERIFICATION_NOTES.md.
%  The original draft of this chapter presented specific pilot-study
%  RESULTS (case counts, Cohen's kappa values, coder agreement) for a
%  pilot that has never run and has not even been submitted for ethics
%  approval. That content has been removed entirely, not edited down --
%  it described data that does not exist. This chapter now describes the
%  real evaluation methodology this project actually used.
% ============================================================
\chapter{Evaluation Methodology}
\label{chap:evaluation}

This chapter describes the evaluation methodology actually used in this
project. Three real evaluation efforts are reported: an external benchmark
for ISCO-08 (WISCO), a disclosed synthetic benchmark for ISIC Rev.~4 and
ISCED-F~2013 (built because no real labelled dataset exists for either), and
a smaller, secondary synthetic-corpus comparison predating the main
evaluation harness. A fourth, planned evaluation --- the randomised pilot
study --- is described as a protocol only: \textbf{it has not been executed,
no data from it exists, and it had not been submitted for ethics approval as
of the writing of this thesis.}

% ------------------------------------------------------------
\section{ISCO-08 Evaluation: The WISCO External Benchmark}
\label{sec:eval-wisco}
% ------------------------------------------------------------

WISCO (World database of ISCO Occupations) is a real, external,
occupation-coding dataset used as this project's primary ISCO-08 evaluation
source, precisely because it is independent of this project's own catalogue
construction. The full, official heldout split used for this thesis's
headline results contains \textbf{18,747 cases}. Two retrieval
configurations were evaluated on the identical case set: flat retrieval
(single-collection direct search) and hierarchical retrieval (four-stage,
parent-filtered beam search) --- see Chapter~\ref{chap:architecture} for
both pipelines' design. Results are reported exactly as measured in
Chapter~\ref{chap:results}, including a negative finding (hierarchical
underperforming flat) that an earlier internal draft of this thesis did not
report.

Beyond the baseline configuration, several additive experiments were run on
the same 18,747-case heldout set to isolate the effect of specific changes:
a larger embedding model (\texttt{multilingual-e5-large} vs.\ the
production-default \texttt{-e5-small}), real official-text catalogue
enrichment (replacing thin \texttt{"\{code\} \{title\}"} embedding text with
the official ILO source document's actual definitions and examples), and
the combination of both. Each experiment's exact result, including whether
LLM re-ranking was tested with it, is reported in
Chapter~\ref{chap:results} --- not summarised here to avoid restating
numbers in two places that could drift out of sync.

A separate, smaller check evaluated whether LLM-based re-ranking of
retrieval candidates improves accuracy, repeated independently across
different retrieval configurations and providers (Groq, and separately a
local Ollama model, since the project's Anthropic account had zero usable
credit throughout evaluation --- Claude results were never used to report
this check, and are not claimed here as having been obtained). The
consistent finding across repeated, independent checks is reported plainly
in Chapter~\ref{chap:results}, whatever it turned out to be.

% ------------------------------------------------------------
\section{ISIC Rev.~4 / ISCED-F 2013 Evaluation: A Disclosed Synthetic Benchmark}
\label{sec:eval-synthetic}
% ------------------------------------------------------------

WISCO is occupation-only --- it carries no industry or field-of-study gold
labels, so it cannot evaluate the ISIC Rev.~4 or ISCED-F~2013 classifiers.
No other real, labelled, multilingual dataset for these two standards was
found; the IPUMS International project was contacted directly to ask
whether it held original-language verbatim occupation/industry/education
text for its Arabic, Hindi, and Urdu samples specifically, and confirmed in
writing that it does not have access to, and could not distribute, any such
data (full correspondence transcript in
\texttt{Documentation/Phase\_2/Week\_1/ipums\_correspondence\_log.md}).

Given this, a synthetic evaluation dataset was built using a real, standard
technique: \textbf{taxonomy-grounded LLM paraphrase data augmentation}. An
LLM is given a classification category's own official definition (sourced
from the real ILO/UNESCO structure documents, the same data used to build
this project's catalogue-enrichment work) and instructed to produce a
casual, first-person sentence a survey respondent might plausibly say,
explicitly avoiding the official terminology so a keyword-matching
classifier cannot trivially win by echoing the source text back. The gold
label is fixed by construction --- it is exactly the category whose
definition generated the text --- not inferred or scored by a second model.

\textbf{This measures a specific, narrower thing than real-respondent
evaluation, stated precisely rather than implied to be equivalent:} it
measures how well a method recovers the class its own generating prompt was
built from, not how the method would perform on genuine survey responses.
It is the best available signal in the absence of real data, useful for
regression testing and directional comparison, but it is never cited in
this thesis as WISCO-equivalent or pilot-grade evidence.

\textbf{Real, disclosed limitations of this dataset, stated exactly:}
generation used a commercial LLM API (Groq) subject to a real daily
token quota; coverage of the full target (all matched ISIC/ISCED-F codes
$\times$ all 6 language codes) was therefore built up incrementally across
multiple sessions rather than completed in one pass. The exact coverage
figure and per-language breakdown achieved is reported in
Chapter~\ref{chap:results}, together with a real data-quality defect this
project found and fixed during construction: a small number of generation
calls returned a plain LLM safety-filter refusal (e.g.\ for a category whose
official examples include sensitive terms such as escort/dating services)
which an early version of the generation pipeline did not detect, silently
storing the refusal text as if it were valid respondent data. An automated,
independent quality-check script was built specifically to catch this and
related defects (script contamination, wrong-script generation, verbatim
title leakage) across every row, not just the affected ones, and the
affected rows were regenerated and re-verified before being used in any
reported result.

A companion evaluation used the Gulf-dialect portion of this synthetic
dataset to test the real, hand-built 79-rule Gulf-Arabic normalisation
dictionary against genuinely dialectal (not Modern Standard Arabic)
generated text --- the first time this project had any dialectal Arabic
data to test that component against at all. The result is reported in
Chapter~\ref{chap:results}.

% ------------------------------------------------------------
\section{A Secondary, Pre-Existing Synthetic-Corpus Comparison}
\label{sec:eval-legacy}
% ------------------------------------------------------------

A separate, smaller framework (\texttt{eval/legacy\_thesis\_ch6/evaluate.py},
predating the main evaluation harness above) compares BM25 sparse
retrieval, flat dense retrieval, and hierarchical dense retrieval on a
100-item \emph{synthetic} corpus.
\textbf{The original draft's gold-labelling claim was checked directly and
is false.} Its ``two ILO-certified coders; adjudicated by senior coder''
claim does not match the real script: the gold labels are a hardcoded
Python list, \texttt{TEST\_CASES: list[tuple[str, str]]} (job description,
true ISCO 4-digit code), 100 English cases plus a separate 30-case Arabic
set (\texttt{ARABIC\_TEST\_CASES}) --- no evidence of external human coding
or adjudication exists in the script, and no separate Hindi/Urdu/Tagalog
case sets were found, contradicting the original draft's claimed
``equal distribution across the 5 languages.'' This is disclosed plainly:
this secondary corpus's gold labels are developer-authored, not
professionally coded.
This framework is real and still in use (its BM25 baseline class is
imported by the main evaluation harness), but its numbers are
\textbf{never} the primary, citable result for this thesis --- the WISCO
result (Section~\ref{sec:eval-wisco}) is. Where this secondary comparison's
results are reported at all (Chapter~\ref{chap:results}), they are labelled
explicitly as synthetic-corpus, non-primary results.

% ------------------------------------------------------------
\section{Planned Pilot Study --- Protocol Only, Not Executed}
\label{sec:eval-pilot}
% ------------------------------------------------------------

\textbf{No pilot data exists. This section describes a plan, not a
completed evaluation.} A draft protocol
(\texttt{Documentation/Phase\_2/Week\_1/module\_e\_pilot\_protocol\_draft.md})
exists, itself marked draft/not-submitted/not-reviewed, built from
already-decided project parameters and standard human-subjects protocol
structure, with every institution-specific detail (ethics committee
calendar, exact recruitment criteria, translated consent materials) left as
an open item pending the supervisor and the institute's ethics review body.

The intended design, stated as a plan: a parallel-arm randomised comparison
($n=30$, 15 AI-administered, 15 traditional-interviewer-administered),
measuring completion time, item nonresponse, internal contradiction rate
(as flagged by the Semantic Relation Engine or Validation Agent),
respondent satisfaction, and per-interview cost. Given $n=30$, this is
explicitly scoped as a feasibility pilot, not a definitively powered
comparison --- any eventual analysis should state this limitation plainly,
not present pilot-scale results as conclusive.

\textbf{An earlier internal draft of this thesis stated specific pilot
results} --- a participant sample description, a count of ISCO-08
classification events, and a Cohen's $\kappa$ inter-coder-agreement value
for double-coded pilot data. \textbf{None of this exists.} It has been
removed from this thesis entirely rather than edited, because it described
data from a study that has never run, not a stale or approximate figure
from one that had.

% ------------------------------------------------------------
\section{Statistical Methods}
\label{sec:eval-stats}
% ------------------------------------------------------------

For the real evaluations reported in Chapter~\ref{chap:results}: paired
accuracy comparisons use McNemar's exact test; proportions are reported
with Wilson 95\% confidence intervals. Both are computed directly (this
project's evaluation code has no \texttt{scipy} dependency; these are
hand-implemented against the standard formulas) rather than taken from an
external library without verification.

For the pilot study, \emph{if and when it runs}: the planned approach is
Welch's $t$-test for completion time (with Cohen's $d$), Mann--Whitney $U$
for CSAT (with rank-biserial $r$), and Fisher's exact test for
nonresponse/contradiction-rate comparisons --- standard choices for a
small-$n$, two-arm design, stated here as a plan to be followed, not as
having already been applied to real data.
